\documentclass{article}
\usepackage{amsmath,amssymb}
\begin{document}

\section*{Higham Chapter 6, Problem 6.8: Lean-Friendly Proof Route}

I am formalizing a source-facing version of the following standard theorem.
For every complex square matrix \(A\) and every \(\delta>0\), there is a
consistent/subordinate matrix norm, depending on \(A\) and \(\delta\), such
that
\[
  \|A\| \le \rho(A)+\delta,
\]
where \(\rho(A)\) is the spectral radius.  Consequently, if \(\rho(A)<1\),
there is a consistent/subordinate matrix norm for which \(\|A\|<1\).

\section*{Local formalization surface}

The local Lean development has the following source-facing concepts.

\begin{itemize}
\item \(C^n\) is represented as functions \(\mathrm{Fin}\ n \to \mathbb C\).
\item A vector norm is a function \(\nu:C^n\to\mathbb R\) satisfying
  nonnegativity, definiteness, complex homogeneity, and the triangle
  inequality.
\item A mixed subordinate upper bound is the predicate
  \[
    \forall x,\quad \nu(Tx)\le C\nu(x).
  \]
\item A mixed subordinate norm value is the least real \(C\) with that
  bound.
\item Already proved: any eigenvalue modulus, and hence any supplied maximum
  eigenvalue modulus \(\rho\), is bounded above by any subordinate norm value
  induced by the same vector norm.
\item There is a generic existence theorem: once a vector norm \(\nu\) and a
  checked bound \(\nu(Ax)\le C\nu(x)\) are available in nonzero dimension,
  Lean can produce a least subordinate norm value \(c\le C\).
\end{itemize}

\section*{What I need}

Please give a rigorous proof route optimized for Lean.

I especially need advice on the smallest mathematically honest theorem to
formalize first if the library does not expose a ready-made Schur/Jordan
triangularization theorem.

Candidate route:
\begin{enumerate}
\item Assume or construct a similarity transform putting \(A\) into upper
  triangular form \(T=S^{-1}AS\), with diagonal entries \(\lambda_i\) and
  \(|\lambda_i|\le\rho\).
\item For a positive scale parameter \(r\), define a weighted sup norm in the
  triangular coordinates, for example
  \[
    \nu_r(x)=\max_i r^i |(S^{-1}x)_i|
  \]
  or an equivalent convention.
\item Show that the subordinate infinity norm of \(T\) after diagonal scaling
  is bounded by
  \[
    \max_i |\lambda_i| + \max_i\sum_{j>i} |T_{ij}|\,r^{j-i}
  \]
  or by the corresponding row-sum formula.
\item Use finiteness to choose \(r>0\) small enough that all off-diagonal row
  sums are at most \(\delta\), giving \(\nu_r(Ax)\le(\rho+\delta)\nu_r(x)\).
\item Then use the generic least-bound theorem to obtain a subordinate norm
  value \(c\le\rho+\delta\).
\item For the consequence, if \(\rho<1\), choose \(\delta=(1-\rho)/2\).
\end{enumerate}

\section*{Questions}

\begin{enumerate}
\item Is the candidate weighted-sup route correct as stated?  If not, fix the
  scaling convention and inequality.
\item What is the cleanest finite-dimensional real inequality for choosing the
  scale parameter?  I want something easy to prove in Lean over finite
  sums/maxima.
\item If full triangularization is too large, what conditional theorem should
  I formalize so it is a genuine dependency rather than a disguised restatement
  of the target theorem?
\item Which hidden assumptions should be explicit?  For example: \(0<n\),
  \(\delta>0\), supplied maximum spectral modulus, invertibility of the change
  of coordinates, diagonal entries bounded by \(\rho\), and strict positivity
  of the scaling parameter.
\item Please give theorem statements and a proof order in Lean-friendly
  language, including the key algebraic identities needed to transfer the
  triangular-coordinate norm back to the original matrix \(A\).
\end{enumerate}

\end{document}
